\chapter{Referensi Cepat Python \& PyTorch}
\label{chap:python_ref}

Appendix ini berfungsi sebagai referensi cepat untuk sintaks dan fungsi penting dalam Python, NumPy, Pandas, dan PyTorch yang sering digunakan dalam pengembangan AI.

\section{Python Dasar untuk Data Science}

\subsection{List Comprehension}
Cara ringkas membuat list.
\begin{lstlisting}[language=Python]
numbers = [1, 2, 3, 4, 5]
squared = [x**2 for x in numbers if x > 2]
# Result: [9, 16, 25]
\end{lstlisting}

\subsection{Lambda Functions}
Fungsi anonim satu baris.
\begin{lstlisting}[language=Python]
add = lambda x, y: x + y
print(add(5, 3)) # Output: 8

# Sering digunakan dengan map/filter
items = [1, 2, 3, 4]
evens = list(filter(lambda x: x % 2 == 0, items))
\end{lstlisting}

\subsection{Dictionary Comprehension}
\begin{lstlisting}[language=Python]
keys = ['a', 'b', 'c']
values = [1, 2, 3]
my_dict = {k:v for k,v in zip(keys, values)}
# {'a': 1, 'b': 2, 'c': 3}
\end{lstlisting}

\subsection{args dan kwargs}
Menangani argumen variabel.
\begin{lstlisting}[language=Python]
def func(*args, **kwargs):
    print(args)   # Tuple dari argumen posisi
    print(kwargs) # Dictionary dari argumen keyword

func(1, 2, name="Alice", age=30)
# (1, 2)
# {'name': 'Alice', 'age': 30}
\end{lstlisting}

\section{NumPy Cheat Sheet}
NumPy adalah fondasi komputasi numerik di Python.

\begin{lstlisting}[language=Python]
import numpy as np

# Membuat Array
a = np.array([1, 2, 3])
b = np.zeros((3, 3))       # Matriks 3x3 isi 0
c = np.ones((2, 4))        # Matriks 2x4 isi 1
d = np.eye(3)              # Matriks identitas 3x3
e = np.random.rand(3, 2)   # Acak uniform [0, 1)
f = np.random.randn(3, 2)  # Acak normal (Gaussian)

# Inspeksi Array
print(a.shape)   # Dimensi array
print(a.ndim)    # Jumlah dimensi
print(a.dtype)   # Tipe data elemen

# Operasi Matematika
x = np.array([1, 2])
y = np.array([3, 4])

print(x + y)     # [4, 6]
print(x * y)     # [3, 8] (Element-wise)
print(x @ y)     # Dot product (1*3 + 2*4 = 11)
print(np.dot(x, y)) # Sama dengan @

# Broadcasting
# Menambahkan vektor ke matriks
m = np.array([[1, 2, 3], [4, 5, 6]])
v = np.array([1, 0, 1])
print(m + v)
# [[2, 2, 4],
#  [5, 5, 7]]

# Reshaping
z = np.arange(12)  # [0, 1, ..., 11]
z_reshaped = z.reshape(3, 4) # Ubah ke 3 baris, 4 kolom

# Indexing & Slicing
print(z_reshaped[0, :])  # Baris pertama, semua kolom
print(z_reshaped[:, 1])  # Semua baris, kolom kedua
print(z_reshaped[z_reshaped > 5]) # Boolean indexing (nilai > 5)
\end{lstlisting}

\section{Pandas Cheat Sheet}
Untuk manipulasi data tabular.

\begin{lstlisting}[language=Python]
import pandas as pd

# Membuat DataFrame
data = {'Name': ['Anna', 'Bob'], 'Age': [24, 30]}
df = pd.DataFrame(data)

# Membaca Data
df = pd.read_csv('data.csv')
df = pd.read_excel('data.xlsx')

# Inspeksi Data
df.head()        # 5 baris pertama
df.info()        # Info tipe data dan null values
df.describe()    # Statistik deskriptif (mean, std, min, max)
df.columns       # Nama kolom
df.shape         # (baris, kolom)

# Seleksi Data
col = df['Name']             # Pilih satu kolom (Series)
cols = df[['Name', 'Age']]   # Pilih beberapa kolom (DataFrame)
row = df.loc[0]              # Pilih baris berdasarkan label index
row_pos = df.iloc[0]         # Pilih baris berdasarkan posisi integer

# Filter Data
adults = df[df['Age'] > 18]
complex_filter = df[(df['Age'] > 18) & (df['Name'] == 'Bob')]

# Manipulasi Data
df['Age_Squared'] = df['Age'] ** 2       # Buat kolom baru
df.drop('Age_Squared', axis=1, inplace=True) # Hapus kolom
df.rename(columns={'Name': 'Full Name'}, inplace=True) # Ganti nama kolom

# Missing Values
df.isnull().sum()            # Hitung nilai null per kolom
df.dropna()                  # Hapus baris dengan nilai null
df.fillna(0)                 # Isi nilai null dengan 0
df['Age'].fillna(df['Age'].mean(), inplace=True) # Imputasi rata-rata

# Grouping & Aggregation
# SELECT Name, AVG(Age) FROM df GROUP BY Name
df.groupby('Name')['Age'].mean()

# Sorting
df.sort_values(by='Age', ascending=False)
\end{lstlisting}

\section{Matplotlib & Seaborn}
Visualisasi data.

\begin{lstlisting}[language=Python]
import matplotlib.pyplot as plt
import seaborn as sns

x = np.linspace(0, 10, 100)
y = np.sin(x)

# Plot Sederhana
plt.figure(figsize=(10, 6))
plt.plot(x, y, label='Sin Wave', color='blue', linestyle='--')
plt.title('Grafik Sinus')
plt.xlabel('X-axis')
plt.ylabel('Y-axis')
plt.legend()
plt.grid(True)
plt.show()

# Scatter Plot
plt.scatter(x, y, marker='x')

# Histogram
data = np.random.randn(1000)
plt.hist(data, bins=30, alpha=0.7)

# Seaborn Heatmap (Korelasi)
corr = df.corr()
sns.heatmap(corr, annot=True, cmap='coolwarm')
\end{lstlisting}

\section{PyTorch Cheat Sheet}
Library Deep Learning utama.

\subsection{Tensor Basics}
\begin{lstlisting}[language=Python]
import torch

# Membuat Tensor
x = torch.tensor([1, 2, 3])
y = torch.ones(2, 3)
z = torch.rand(2, 3)

# Pindah ke GPU
device = torch.device('cuda' if torch.cuda.is_available() else 'cpu')
x = x.to(device)

# Konversi NumPy <-> Tensor
np_arr = x.cpu().numpy()
tensor_from_np = torch.from_numpy(np_arr)

# Operasi (mirip NumPy)
a = torch.randn(3, 3)
b = torch.randn(3, 3)
c = a + b
d = torch.matmul(a, b) # Matrix multiplication
\end{lstlisting}

\subsection{Autograd (Automatic Differentiation)}
\begin{lstlisting}[language=Python]
# requires_grad=True memberitahu PyTorch untuk melacak operasi
w = torch.tensor([1.0], requires_grad=True)
x = torch.tensor([2.0])
b = torch.tensor([3.0])

# Forward pass
y = w * x + b  # y = 1*2 + 3 = 5

# Backward pass (hitung gradien)
y.backward()

# Gradien dy/dw = x = 2.0
print(w.grad)
\end{lstlisting}

\subsection{Neural Network Modules (torch.nn)}
\begin{lstlisting}[language=Python]
import torch.nn as nn
import torch.nn.functional as F

class SimpleNet(nn.Module):
    def __init__(self):
        super(SimpleNet, self).__init__()
        # Definisi layer
        self.fc1 = nn.Linear(in_features=784, out_features=128)
        self.fc2 = nn.Linear(128, 64)
        self.fc3 = nn.Linear(64, 10)
        self.dropout = nn.Dropout(p=0.2)

    def forward(self, x):
        # Definisi aliran data
        x = x.view(-1, 784) # Flatten input
        x = F.relu(self.fc1(x))
        x = self.dropout(x)
        x = F.relu(self.fc2(x))
        x = self.fc3(x) # Output raw logits (tanpa softmax jika pakai CrossEntropyLoss)
        return x

model = SimpleNet().to(device)
\end{lstlisting}

\subsection{Training Loop Template}
\begin{lstlisting}[language=Python]
import torch.optim as optim

# Hyperparameters
lr = 0.001
epochs = 5

# Loss dan Optimizer
criterion = nn.CrossEntropyLoss()
optimizer = optim.Adam(model.parameters(), lr=lr)

# Training Loop
for epoch in range(epochs):
    for batch_idx, (data, targets) in enumerate(train_loader):
        data = data.to(device)
        targets = targets.to(device)

        # 1. Forward Pass
        scores = model(data)
        loss = criterion(scores, targets)

        # 2. Backward Pass
        optimizer.zero_grad() # Reset gradien ke 0
        loss.backward()       # Hitung gradien baru

        # 3. Update Weights
        optimizer.step()

    print(f"Epoch {epoch} Loss: {loss.item()}")
\end{lstlisting}

\subsection{Saving & Loading Models}
\begin{lstlisting}[language=Python]
# Save seluruh model (arsitektur + bobot)
torch.save(model, 'model.pth')

# Load seluruh model
model = torch.load('model.pth')

# (Recommended) Save hanya state dict (bobot)
torch.save(model.state_dict(), 'weights.pth')

# Load weights (harus instansiasi model dulu)
model = SimpleNet()
model.load_state_dict(torch.load('weights.pth'))
model.eval() # Set ke mode evaluasi (matikan dropout/batchnorm)
\end{lstlisting}

\subsection{Dataset & DataLoader}
\begin{lstlisting}[language=Python]
from torch.utils.data import Dataset, DataLoader

class CustomDataset(Dataset):
    def __init__(self, csv_file):
        self.data = pd.read_csv(csv_file)

    def __len__(self):
        return len(self.data)

    def __getitem__(self, idx):
        # Ambil satu sampel
        sample = self.data.iloc[idx]
        features = torch.tensor(sample[:-1], dtype=torch.float32)
        label = torch.tensor(sample[-1], dtype=torch.long)
        return features, label

dataset = CustomDataset('train.csv')
train_loader = DataLoader(dataset, batch_size=32, shuffle=True)
\end{lstlisting}

\section{Common Errors & Solusi}

\begin{itemize}
    \item \textbf{RuntimeError: Expected object of scalar type Long but got Float}:
    Biasanya terjadi pada target label untuk CrossEntropyLoss. Target harus berupa integer (Long), bukan float.
    \item \textbf{RuntimeError: Input type (torch.FloatTensor) and weight type (torch.cuda.FloatTensor) should be the same}:
    Anda lupa memindahkan input data ke GPU (.to(device)), sementara model sudah di GPU.
    \item \textbf{CUDA out of memory}:
    Batch size terlalu besar. Kurangi batch size atau gunakan \texttt{with torch.no\_grad():} saat evaluasi.
    \item \textbf{Gradient is None}:
    Anda mungkin lupa memanggil \texttt{loss.backward()}, atau tensor tidak memiliki \texttt{requires\_grad=True}.
\end{itemize}
